\originalsection{磁力的方向、选择速度与轨迹}\label{tb:selector}
磁场可由运动电荷所受横向力测量: $B=F/(|q|v\sin\theta)$, 因而T等于$\mathrm{N/(A\,m)}$. 速度与场平行时力为零, 不能用这一次零力测量认定磁场不存在. 同一场中的负电荷须把正电荷右手叉积再反向.

速度选择器让$\ee=E\uv y$, $\bb=B\uv z$, 入射速度$v\uv x$. 以$q>0$说明方向, 电场力向$+y$, 磁力$qv\uv x\times B\uv z=-qvB\uv y$, 因而无偏转条件$v=E/B$. 若电场、磁场和速度不互相垂直, 一个标量比值不能代替完整矢量方程. 负电荷两力方向同时反转, 相消条件仍是同一速率. 选择器挑出速度, 后续只有磁场的分析区再以曲率给$m/|q|=BR/v$.

\begin{example}
\textbf{原8.9.3: 同压加速的两种粒子.} 
正粒子A电荷$q$, B电荷$2q$, 从静止经过同一加速势差$\Delta V$后进入垂直磁场. 半圆轨道半径分别为$R,2R$. 求$m_A/m_B$.
\end{example}
\begin{solution}
$\tfrac12mv^2=q\Delta V$, $R=mv/(qB)$, 合并$R^2=2m\Delta V/(qB^2)$. 所以
\[\frac{R_A^2}{R_B^2}=\frac{m_A/q}{m_B/(2q)}=\frac14,\qquad \frac{m_A}{m_B}=\frac18.\]
若只是相同速度而非相同加速电压, 则质量比会不同. 题中的“从静止”与“同势差”不可省略.
\end{solution}

把旋转写成实际时间函数还可核查符号. $q>0$, $\bb=B\uv z$, 初始$\vct v=v_\perp\uv x+v_\parallel\uv z$, 则
\[v_x=v_\perp\cos\omega_ct,\quad v_y=-v_\perp\sin\omega_ct,\quad z=z_0+v_\parallel t.\]
从$+z$看轨迹顺时针, 半径固定. 电场沿$z$同时存在时$v_z$随时间改变, 轨迹的螺距也不再固定. 非均匀磁场可能使半径随位置变化, 本章的等速圆公式不能直接覆盖任意轨迹.

\originalsection{曲线导线与机械能}\label{tb:rod-magnetic}
\begin{example}
\textbf{原8.1: 闭合半圆导线.} 
半径$R$的上半圆和直径构成闭合回路, 电流沿直径从$(-R,0)$到$(R,0)$, 再沿上半圆返回, 均匀磁场$B\uv y$. 求弧和直径各自受力.
\end{example}
\begin{solution}
均匀场可移到积分外, $\vct F_{\rm arc}=I[\rr_{\rm end}-\rr_{\rm start}]\times\bb=-2IRB\uv z$, 直径受$+2IRB\uv z$, 合力零. 有限曲线上局部力方向变化, 总力取决于端点位移, 不是弧长乘$I B$. 该结论只限均匀$\bb$.
\end{solution}

\begin{example}
\textbf{原8.9.1: 无滑动滚动的导棒.} 
均匀实心圆柱棒质量$m$, 半径$R$, 架在间距$\ell$、可行长度$a$的平行导轨上. 电流保持$I$, 磁場与轨道平面垂直, 棒从静止无滑动滚过距离$a$. 求离开时速度.
\end{example}
\begin{solution}
推力$F=I\ell B$, 功$Fa$. 轴向转动惯量记$\mathcal I=mR^2/2$, 避免与电流$I$混淆. 无滑动约束$\omega=v/R$, 静摩擦在理想固定接触点不作功但提供力矩, 因而
\[I\ell Ba=\tfrac12mv^2+\tfrac12\mathcal I(v/R)^2=\tfrac34mv^2,
\quad v=\sqrt{\frac{4I\ell Ba}{3m}}.\]
电源维持电流时提供能量, 并要对抗棒运动的感应电动势. 把它称为“磁力对每个载流子作功”会与$\vct F_B\cdot\vct v=0$矛盾; 完整系统通过电场、约束与电源交换能量. 原文一处把无滑动误写成with slipping, 此处按实际$\omega=v/R$条件修正.
\end{solution}

\begin{example}
\textbf{原8.9.2、8.9.4: 悬杆与环.} 
(a)线密度$\lambda_m$的水平杆, 外场垂直纸面向外, 要求悬挂绳张力为零. (b)水平环半径$r$, 棒磁铁N端在环下方, 环处场与竖直夹角$\theta$, 场的径向分量向外. 电流从上看顺时针, 求环受力.
\end{example}
\begin{solution}
(a)向左电流给向上磁力, $I\ell B=\lambda_m\ell g$, 所以$I=\lambda_mg/B$. 这是竖直力平衡, 并不证明横向稳定.

(b)$\dd\vct\ell=-\uv\phi\dd s$, $\bb=B\sin\theta\uv r+B\cos\theta\uv z$. 叉积为$I B\sin\theta\dd s\uv z-I B\cos\theta\dd s\uv r$. 径向力环积分相消, 总力
\[\vct F=2\pi rIB\sin\theta\uv z.\]
若反转电流方向力反转; 总力来自径向场分量, 只以$B_z$代入会遗漏它.
\end{solution}

\originalsection{磁化的等效电流、地磁与磁滞}\label{tb:magnetization}
许多小电流环组成的物质, 相邻环边电流在内部抵消, 只有边界或非均匀处留下等效电流. 定义磁化强度$\vct M$为单位体积磁矩, 单位A/m. 均匀磁化圆柱若$\vct M=M\uv z$, 侧面等效面电流$\vct K_b=\vct M\times\uv n=M\uv\phi$, 类似螺线管; 一般$\jj_b=\nabla\times\vct M$. 安培环量计自由加束缚电流, 因而定义$\vct H=\bb/\mu_0-\vct M$可写$\oint\vct H\cdot\dd\vct\ell=I_{f,\rm enc}$.

弱顺磁材料的既有微观磁矩有部分沿外场取向, 热运动妨碍排列, 通常$\chi_m>0$且温度会影响大小. 抗磁响应诱导反向磁矩, 通常$\chi_m<0$. 铁磁材料有磁畴, 增大外场使畴壁移动和磁矩转向, 最终饱和; 降低外场时会保留剩磁, 需反向矫顽场才能消去. 因而$M(H)$不是同一条直线, 不能用一个恒定$\mu$覆盖磁滞回线. 回线面积反映每循环的能量损耗, 也解释变压器磁芯发热的一部分来源.

地球远处近似偶极场, 但磁矩朝向使地理北方附近的外部场具有磁性南极特征. 当地磁场同时含水平和竖直分量, 指南针沿水平投影定向, 悬在三维自由轴上的磁针还会显示倾角. 原9.5.1列出MIT所在地的历史观测值, 本稿把它作为教材当时的案例, 不把该值当作2026年当前导航数据. 偶极公式可由纬度解释场的数量级和倾斜, 真正地磁含非偶极分量并随时间变化.
